%%% Task Manager — Yazılım Kalite Güvencesi
%%% Aşama 2: İlk Statik Analiz Raporu
%%% IEEE Conference / IEEEtran formatı
%%% Overleaf'te "New Project -> Blank Project" -> bu dosyayı yapıştır ve derle.
%%% Derleyici: pdfLaTeX (Overleaf -> Menu -> Compiler: pdfLaTeX).

\documentclass[conference]{IEEEtran}
\IEEEoverridecommandlockouts

% --- Türkçe karakter ve font desteği ----------------------------------------
\usepackage[utf8]{inputenc}
\usepackage[T1]{fontenc}
\usepackage[turkish]{babel}
\usepackage{lmodern}

% --- Tablo, listeleme ve dizgi paketleri ------------------------------------
\usepackage{cite}
\usepackage{amsmath,amssymb,amsfonts}
\usepackage{graphicx}
\usepackage{textcomp}
\usepackage{xcolor}
\usepackage{booktabs}
\usepackage{array}
\usepackage{hyperref}
\usepackage{url}
\usepackage{caption}

\hypersetup{
  colorlinks=true,
  linkcolor=black,
  urlcolor=blue,
  citecolor=black
}

\def\BibTeX{{\rm B\kern-.05em{\sc i\kern-.025em b}\kern-.08em
    T\kern-.1667em\lower.7ex\hbox{E}\kern-.125emX}}

\begin{document}

% ----------------------------------------------------------------------------
\title{Task Manager Uygulaması Üzerinde\\
İlk Statik Analiz Raporu (Aşama 2)%
\thanks{Yazılım Kalite Güvencesi Dersi Dönem Projesi Aşama~2 Teslimi.}}

\author{\IEEEauthorblockN{Furkan Öztürk}
\IEEEauthorblockA{\textit{Yazılım Mühendisliği Bölümü} \\
\textit{Öğrenci No: 230229083}\\
2026-04-22}}

\maketitle

% ----------------------------------------------------------------------------
\begin{abstract}
Bu çalışma, Yazılım Kalite Güvencesi dönem projesinin ikinci
aşaması kapsamında geliştirilen \emph{Task Manager} web uygulamasının ilk
statik analizini sunmaktadır. Uygulama; \texttt{Node.js}, \texttt{Express}
tabanlı bir REST API sunucusu ve \texttt{React/TypeScript} tabanlı bir
istemci uygulamasından oluşmaktadır. Analiz, endüstride yaygın kullanılan
\textbf{SonarQube~10.6 Community} ve \textbf{ESLint~8.57} araçları ile,
yalnızca proje kaynak kodu üzerinde gerçekleştirilmiştir. Sonuç olarak
5153~NCLOC üzerinden SonarQube~83 bulgu (14~BUG ve 69~CODE\_SMELL) ve
ESLint~41 bulgu raporlamıştır. Güvenlik tarafı stabil
(\texttt{vulnerabilities}=0, \texttt{security\_rating}=A) iken
güvenilirlik derecelendirmesi \emph{D} seviyesindedir. Aşama~3 kod
güncelleme sürecinde önceliklendirmek üzere somut aksiyon planı
oluşturulmuştur.
\end{abstract}

\begin{IEEEkeywords}
yazılım kalite güvencesi, statik analiz, SonarQube, ESLint, teknik borç,
TypeScript, React
\end{IEEEkeywords}

% ----------------------------------------------------------------------------
\section{Giriş}

Yazılım kalite güvencesinin temel pratiklerinden biri, kod tabanı üzerinde
periyodik olarak statik analiz araçlarını çalıştırarak ölçülebilir
metrikler elde etmek ve bu metrikler üzerinden iyileştirme döngülerini
yönetmektir~\cite{b1,b2}. Bu dönem projesi dört aşamalı bir
süreç olarak planlanmış olup bu rapor \emph{Aşama~2 — İlk Statik Analiz}
çıktısıdır. Aşama~2'de elde edilen taban çizgisi (\emph{baseline}),
Aşama~3'te yapılacak kod güncellemelerinin ardından Aşama~4'te aynı
komut ve ayarlar ile tekrarlanacak ve iki ölçüm karşılaştırılacaktır.

\section{Kullanılan Araçlar ve Konfigürasyon}

\subsection{Çalışma Ortamı}

\begin{table}[htbp]
\caption{Çalışma ortamı bileşenleri}
\label{tab:env}
\centering
\begin{tabular}{@{}ll@{}}
\toprule
\textbf{Bileşen} & \textbf{Versiyon} \\
\midrule
İşletim sistemi & Windows 11 Pro \\
Node.js & 22.15.0 \\
npm & 10.9.2 \\
Java (sonar-scanner için) & Eclipse Adoptium JDK 21.0.8 \\
Docker Engine & 28.3.2 \\
\bottomrule
\end{tabular}
\end{table}

\subsection{ESLint Konfigürasyonu}

ESLint~8.57.0 sürümü, TypeScript ve React için resmi eklentilerle
kullanılmıştır. Kullanılan kural seti Tablo~\ref{tab:eslint-rules}
ve Tablo~\ref{tab:eslint-deps}'de özetlenmiştir.

\begin{table}[htbp]
\caption{ESLint bağımlılıkları}
\label{tab:eslint-deps}
\centering
\begin{tabular}{@{}ll@{}}
\toprule
\textbf{Paket} & \textbf{Versiyon} \\
\midrule
\texttt{eslint} & 8.57.0 \\
\texttt{@typescript-eslint/parser} & 7.18.0 \\
\texttt{@typescript-eslint/eslint-plugin} & 7.18.0 \\
\texttt{eslint-plugin-react} & 7.35.0 \\
\texttt{eslint-plugin-react-hooks} & 4.6.2 \\
\bottomrule
\end{tabular}
\end{table}

\begin{table}[htbp]
\caption{Kullanılan kural setleri (extends)}
\label{tab:eslint-rules}
\centering
\begin{tabular}{@{}ll@{}}
\toprule
\textbf{Kural Seti} & \textbf{Kapsam} \\
\midrule
\texttt{eslint:recommended} & tüm kod tabanı \\
\texttt{plugin:@typescript-eslint/recommended} & tüm kod tabanı \\
\texttt{plugin:react/recommended} & yalnızca \texttt{client/src/} \\
\texttt{plugin:react-hooks/recommended} & yalnızca \texttt{client/src/} \\
\texttt{plugin:react/jsx-runtime} & yalnızca \texttt{client/src/} \\
\bottomrule
\end{tabular}
\end{table}

Ek olarak \texttt{no-var}=\textit{error},
\texttt{prefer-const}=\textit{warn},
\texttt{eqeqeq}=\textit{smart/warn},
\texttt{@typescript-eslint/no-unused-vars}=\textit{warn}
(\texttt{\_} ön ekli olanlar hariç) ve
\texttt{@typescript-eslint/no-explicit-any}=\textit{warn} olarak
tanımlanmıştır. Test dosyalarında \texttt{no-explicit-any} kuralı
bilinçli olarak devre dışı bırakılmıştır.

\textbf{Kapsam:} \texttt{src/}, \texttt{client/src/}, \texttt{tests/}
(\texttt{.ts}, \texttt{.tsx} uzantıları). \texttt{node\_modules/},
\texttt{dist/}, \texttt{build/}, \texttt{coverage/}, \texttt{client/dist/},
\texttt{data/}, \texttt{reports/} ve config dosyaları hariç tutulmuştur.

\textbf{Komut:}
\begin{small}
\begin{verbatim}
npm run lint:stage1
# eslint --ext .ts,.tsx --format json \
#   --output-file reports/stage1/eslint.json \
#   src client/src tests
\end{verbatim}
\end{small}

\subsection{SonarQube Konfigürasyonu}

SonarQube~10.6 Community sürümü, yerel Docker konteynerinde
(\texttt{sonarqube:10.6-community}) çalıştırılmıştır. Tarayıcı olarak
\texttt{sonar-scanner~3.1.0} (npm paketi) ve Java~21 kullanılmıştır.
Kalite profili olarak SonarQube'un yerleşik \textbf{Sonar way} profili
(TypeScript/JavaScript) seçilmiştir.

\texttt{sonar-project.properties} dosyasındaki temel ayarlar:
\begin{itemize}
\item \texttt{sonar.projectKey = task-manager}
\item \texttt{sonar.sources = src,client/src}
\item \texttt{sonar.tests = tests}
\item \texttt{sonar.exclusions = **/node\_modules/**,\
  **/dist/**,**/build/**,**/coverage/**,\
  reports/**,data/**,client/dist/**}
\item \texttt{sonar.typescript.tsconfigPaths =\
  tsconfig.json,client/tsconfig.json}
\end{itemize}

Yönerge gereği BUG ve CODE\_SMELL issue tipleri; BLOCKER, CRITICAL,
MAJOR, MINOR, INFO severity seviyeleri ve on üç adet metrik
raporlanmıştır (bkz. Bölüm~\ref{sec:bulgular}).

Raporlar SonarQube Web API'si aracılığıyla JSON formatında çekilmiştir:
\texttt{/api/issues/search} ile issue listesi ve
\texttt{/api/measures/component} ile metrikler alınmıştır. Bu işlem
projenin \texttt{scripts/fetch-sonar-reports.mjs} script'i ile
otomatikleştirilmiştir.

% ----------------------------------------------------------------------------
\section{Bulgular}
\label{sec:bulgular}

\subsection{ESLint Bulguları}

Toplam 7 dosyada 15 hata (\textit{error}) ve 26 uyarı (\textit{warning}),
yani 41~bulgu raporlanmıştır. Bulguların kural bazında dağılımı
Tablo~\ref{tab:eslint-rules-hit}'de, dosya bazında dağılımı
Tablo~\ref{tab:eslint-files}'de sunulmuştur.

\begin{table}[htbp]
\caption{ESLint — kural bazında bulgu dağılımı}
\label{tab:eslint-rules-hit}
\centering
\begin{tabular}{@{}lrl@{}}
\toprule
\textbf{Kural} & \textbf{Adet} & \textbf{Seviye} \\
\midrule
\texttt{@typescript-eslint/no-explicit-any} & 36 & error/warn \\
\texttt{react/no-unescaped-entities}        &  2 & error \\
\texttt{@typescript-eslint/no-unused-vars}  &  2 & warn \\
\texttt{react-hooks/exhaustive-deps}        &  1 & warn \\
\midrule
\textbf{Toplam} & \textbf{41} & \\
\bottomrule
\end{tabular}
\end{table}

\begin{table}[htbp]
\caption{ESLint — dosya bazında bulgu dağılımı}
\label{tab:eslint-files}
\centering
\begin{tabular}{@{}lrr@{}}
\toprule
\textbf{Dosya} & \textbf{Hata} & \textbf{Uyarı} \\
\midrule
\texttt{client/src/pages/AdminPage.tsx}        &  5 &  0 \\
\texttt{client/src/pages/TasksPage.tsx}        & 10 &  1 \\
\texttt{src/routes/tasks.ts}                   &  0 &  1 \\
\texttt{src/services/projectService.ts}        &  0 & 19 \\
\texttt{src/services/taskService.ts}           &  0 &  3 \\
\texttt{tests/unit/projectService.test.ts}     &  0 &  1 \\
\texttt{tests/unit/taskService.test.ts}        &  0 &  1 \\
\midrule
\textbf{Toplam} & \textbf{15} & \textbf{26} \\
\bottomrule
\end{tabular}
\end{table}

\subsection{SonarQube Bulguları}

SonarQube tarafında toplam 83 issue (14 BUG ve 69 CODE\_SMELL)
raporlanmıştır. Ölçülen metrikler Tablo~\ref{tab:sonar-measures}'de,
severity dağılımı Tablo~\ref{tab:sonar-severity}'de, en çok tetiklenen
kurallar Tablo~\ref{tab:sonar-rules}'de ve en çok bulgu üreten dosyalar
Tablo~\ref{tab:sonar-files}'de sunulmuştur.

\begin{table}[htbp]
\caption{SonarQube metrikleri}
\label{tab:sonar-measures}
\centering
\begin{tabular}{@{}lr@{}}
\toprule
\textbf{Metrik} & \textbf{Değer} \\
\midrule
\texttt{ncloc}                    & 5153 \\
\texttt{complexity}               & 776 \\
\texttt{cognitive\_complexity}    & 2 \\
\texttt{bugs}                     & 14 \\
\texttt{code\_smells}             & 69 \\
\texttt{vulnerabilities}          & 0 \\
\texttt{security\_hotspots}       & 2 \\
\texttt{duplicated\_lines\_density} & \%1.3 \\
\texttt{reliability\_rating}      & 4.0 (D) \\
\texttt{sqale\_rating}            & 1.0 (A) \\
\texttt{security\_rating}         & 1.0 (A) \\
\texttt{sqale\_index}             & 348 dk \\
\texttt{sqale\_debt\_ratio}       & \%0.2 \\
\bottomrule
\end{tabular}
\end{table}

\begin{table}[htbp]
\caption{SonarQube severity dağılımı}
\label{tab:sonar-severity}
\centering
\begin{tabular}{@{}lrrrrrr@{}}
\toprule
\textbf{Tip} & \textbf{BLOCKER} & \textbf{CRITICAL} & \textbf{MAJOR} & \textbf{MINOR} & \textbf{INFO} & \textbf{Top.} \\
\midrule
BUG         & 0 & 1 & 0 & 13 & 0 & 14 \\
CODE\_SMELL & 0 & 1 & 9 & 59 & 0 & 69 \\
\midrule
\textbf{Genel} & 0 & 2 & 9 & 72 & 0 & \textbf{83} \\
\bottomrule
\end{tabular}
\end{table}

\begin{table}[htbp]
\caption{En çok tetiklenen SonarQube kuralları (Top~10)}
\label{tab:sonar-rules}
\centering
\begin{tabular}{@{}rll@{}}
\toprule
\textbf{Adet} & \textbf{Kural} & \textbf{Kısa anlamı} \\
\midrule
25 & \texttt{typescript:S6606} & \texttt{??} yerine \texttt{||} kullanımı \\
15 & \texttt{typescript:S6853} & Etiket / form kontrol bağlaması \\
14 & \texttt{typescript:S6848} & Non-interactive element handler'ı \\
13 & \texttt{typescript:S1082} & DOM event handler tip uyumsuzluğu \\
 6 & \texttt{typescript:S6759} & Eksik \texttt{readonly} React prop'u \\
 4 & \texttt{typescript:S3863} & Yinelenen \texttt{import} \\
 2 & \texttt{typescript:S3358} & İç içe ternary ifade \\
 1 & \texttt{typescript:S2871} & \texttt{sort()} için karşılaştırıcı eksik \\
 1 & \texttt{typescript:S3696} & Literal değer \texttt{throw} edilmiş \\
 1 & \texttt{typescript:S3626} & Gereksiz \texttt{return/break/continue} \\
 1 & \texttt{typescript:S3776} & Cognitive complexity eşiği aşımı \\
\bottomrule
\end{tabular}
\end{table}

\begin{table}[htbp]
\caption{SonarQube — dosya bazında bulgu dağılımı}
\label{tab:sonar-files}
\centering
\begin{tabular}{@{}lr@{}}
\toprule
\textbf{Dosya} & \textbf{Issue} \\
\midrule
\texttt{client/src/pages/TasksPage.tsx}     & 51 \\
\texttt{client/src/pages/AdminPage.tsx}     &  8 \\
\texttt{src/routes/tasks.ts}                &  6 \\
\texttt{client/src/pages/RegisterPage.tsx}  &  5 \\
\texttt{client/src/pages/LoginPage.tsx}     &  4 \\
\texttt{client/src/pages/ReportsPage.tsx}   &  3 \\
\texttt{src/server.ts}                      &  2 \\
Diğer dosyalar (4 adet) & 4 \\
\midrule
\textbf{Toplam} & \textbf{83} \\
\bottomrule
\end{tabular}
\end{table}

\subsubsection{Kritik Bulgular}
İki adet CRITICAL severity bulgu tespit edilmiştir ve her ikisi de
\texttt{src/routes/tasks.ts} dosyasındadır:

\begin{itemize}
\item \texttt{src/routes/tasks.ts:130} (BUG, \texttt{S2871}):
  \texttt{.sort()} çağrısına karşılaştırıcı fonksiyon verilmemiştir;
  elementler alfabetik sıralanacak, sayısal alanlarda yanlış sonuç
  üretecektir.
\item \texttt{src/routes/tasks.ts:113} (CODE\_SMELL, \texttt{S3776}):
  Bir fonksiyonun cognitive complexity değeri 23 olup izin verilen
  eşik 15'tir.
\end{itemize}

\subsubsection{Öne çıkan gözlemler}
\begin{itemize}
\item \texttt{reliability\_rating=D} derecelendirmesinin başlıca sebebi,
  BUG kategorisindeki 1~CRITICAL bulgudur. Bu bulgunun kapatılması
  Aşama~3'ün en öncelikli hedefidir.
\item \texttt{sqale\_rating=A} ve \texttt{sqale\_debt\_ratio=\%0.2}
  değerleri teknik borcun kontrol altında olduğunu göstermektedir
  (toplam 348~dk tahmini borç).
\item \texttt{duplicated\_lines\_density=\%1.3} değeri makuldür; yine de
  Aşama~3'te düşürülmesi hedeflenmektedir.
\item Güvenlik tarafı stabildir: \texttt{vulnerabilities=0},
  \texttt{security\_rating=A}, yalnızca 2 adet güvenlik incelemesi
  gerektiren nokta mevcuttur.
\item \texttt{client/src/pages/TasksPage.tsx} dosyası tek başına
  tüm SonarQube bulgularının yaklaşık \%61'ini (51/83) üretmektedir.
\end{itemize}

% ----------------------------------------------------------------------------
\section{Planlanan Aksiyonlar}

Aşama~3 (Kod Güncelleme) sürecinde aşağıdaki aksiyonlar önceliklendirme
sırasına göre uygulanacaktır.

\subsection{ESLint Kaynaklı Aksiyonlar}
\begin{enumerate}
\item \textbf{\texttt{@typescript-eslint/no-explicit-any} (36 adet):}
  \texttt{any} yerine proje içindeki mevcut tipler, generic'ler veya
  \texttt{unknown}~+~daraltma (\emph{narrowing}) kullanılacaktır.
\item \textbf{\texttt{react/no-unescaped-entities} (2 adet):}
  JSX içindeki \texttt{'} karakterleri \texttt{\&apos;} veya string
  literali olarak kaçırılacaktır.
\item \textbf{\texttt{react-hooks/exhaustive-deps} (1 adet):}
  \texttt{useEffect} bağımlılık dizisi doğru doldurulacak veya
  \texttt{useCallback} ile stabil referans üretilecektir.
\item \textbf{\texttt{@typescript-eslint/no-unused-vars} (2 adet):}
  Kullanılmayan sembol/importlar temizlenecektir.
\end{enumerate}

\subsection{SonarQube Kaynaklı Aksiyonlar}
\begin{enumerate}
\item \textbf{CRITICAL BUG} \texttt{src/routes/tasks.ts:130}:
  \texttt{sort()} çağrısına karşılaştırıcı fonksiyon eklenecektir;
  \texttt{reliability\_rating} D'den A'ya çıkartılmasının önündeki en
  büyük engel budur.
\item \textbf{CRITICAL CODE\_SMELL} \texttt{src/routes/tasks.ts:113}:
  Fonksiyon yardımcı fonksiyonlara bölünerek cognitive complexity
  değeri $\leq 15$ seviyesine düşürülecektir.
\item \textbf{Yoğunluk merkezli refaktör:}
  \texttt{client/src/pages/TasksPage.tsx} küçük bileşenlere bölünecektir.
\item \textbf{Mekanik kural bazlı temizlik:}
  \texttt{S6606} (25), \texttt{S6853} (15), \texttt{S6848} (14),
  \texttt{S1082} (13), \texttt{S6759} (6), \texttt{S3863} (4) kurallarına
  karşılık gelen değişiklikler toplu olarak uygulanacaktır.
\item \textbf{\texttt{security\_hotspots=2}:} SonarQube UI üzerinden
  inceleme yapılıp \emph{Reviewed} statüsüne geçirilecektir.
\end{enumerate}

\subsection{Ölçüm Hedefleri}

Aşama~4 sonunda, Aşama~2'ye kıyasla hedeflenen iyileşmeler
Tablo~\ref{tab:targets}'de verilmiştir.

\begin{table}[htbp]
\caption{Aşama~4 ölçüm hedefleri}
\label{tab:targets}
\centering
\begin{tabular}{@{}lrr@{}}
\toprule
\textbf{Metrik} & \textbf{Aşama 2} & \textbf{Aşama 4 Hedef} \\
\midrule
ESLint error  & 15 & 0 \\
ESLint warning & 26 & $\leq$ 5 \\
SonarQube BUG & 14 & $\leq$ 2 (CRITICAL=0) \\
SonarQube CODE\_SMELL & 69 & $\leq$ 20 \\
CODE\_SMELL (BLOCKER+CRITICAL) & 1 & 0 \\
\texttt{duplicated\_lines\_density} & \%1.3 & $\leq$ \%1.0 \\
\texttt{reliability\_rating} & D (4.0) & A (1.0) \\
\texttt{sqale\_rating} & A (1.0) & A (korunur) \\
\texttt{security\_rating} & A (1.0) & A (korunur) \\
\texttt{sqale\_index} & 348 dk & $\leq$ 120 dk \\
\texttt{security\_hotspots} & 2 & 0 (reviewed) \\
\bottomrule
\end{tabular}
\end{table}

% ----------------------------------------------------------------------------
\section{Sonuç}

Bu rapor, Task Manager uygulamasının Aşama~2 kapsamındaki ilk statik
analiz sonuçlarını sunmuştur. SonarQube ve ESLint çıktıları yönergede
istenen tüm metrik, issue tipi ve severity seviyeleri bakımından tam
olarak raporlanmıştır. Taban çizgisi metrikleri toplam 83~SonarQube ve
41~ESLint bulgusu, \%61'i tek bir dosyada yoğunlaşan bir dağılım ve D
seviyesinde bir \texttt{reliability\_rating} ortaya koymaktadır. Bu
bulgular Aşama~3 kod güncelleme aşamasının önceliklerini netleştirmekte
ve Aşama~4'teki karşılaştırmalı son analizin dayanacağı ölçüm setini
oluşturmaktadır. Tutarlı bir karşılaştırma sağlamak için Aşama~4,
aynı konfigürasyon dosyaları ve aynı \texttt{npm} komutlarıyla
(\texttt{analyze:stage2}) yürütülecektir.

% ----------------------------------------------------------------------------
\begin{thebibliography}{00}
\bibitem{b1} SonarSource, ``SonarQube Documentation,''
  \url{https://docs.sonarsource.com/sonarqube/}, Erişim: 2026-04-22.
\bibitem{b2} ESLint, ``ESLint — Pluggable JavaScript Linter,''
  \url{https://eslint.org/}, Erişim: 2026-04-22.
\bibitem{b3} G. A. Campbell, ``Cognitive Complexity: A New Way of
  Measuring Understandability,'' SonarSource S.A., Teknik Rapor, 2018.
\end{thebibliography}

\end{document}
